% !TEX program = lualatex
% Auto-generated from YAML (v3) — 06: 情報システムの役割と信頼性

\documentclass[aspectratio=43,professionalfonts,handout]{beamer}
\usepackage{beamer_template}

\newcommand{\LectureNum}{06}
\newcommand{\LectureTitle}{情報システムの役割と信頼性}
\newcommand{\CourseName}{情報リテラシー}
\newcommand{\TextPages}{pp.112-117}
\newcommand{\NextTopic}{データベースとデータ管理の基礎}
\newcommand{\NextPages}{pp.118-121}

\begin{document}
% --- Slide 1: title ---

\begin{frame}{\CourseName\quad 第\LectureNum 回「\LectureTitle 」}
  {\small \TermName\quad \TeacherName\quad ／\quad 教科書 \TextPages}

  \vfill

  \begin{block}{今日の流れ}
    \begin{enumerate}
      \item 情報システムのサービスと種類
      \item 情報システムと通貨（電子マネー・仮想通貨）
      \item さまざまな情報システム（集中・分散・移動体）
      \item オープンデータ
      \item 情報システムの信頼性（バックアップ・RAID・稼働率）
    \end{enumerate}
  \end{block}
\end{frame}

% --- 目標 ---

\begin{frame}{今日の目標}
  \begin{block}{この授業が終わったらできること}
    \begin{itemize}
      \item 社会における情報システムの役割と，その信頼性確保の重要性を説明できる。
    \end{itemize}
  \end{block}
\end{frame}

% --- Slide 2: terms ---

\begin{frame}{情報システムの基本用語}
  \begin{description}[labelwidth=6em]
    \item[\kw{情報システム}]
      収集されたデータを利用可能な形に処理し，必要とする場所へ伝達する仕組み。さまざまなハードウェア・ソフトウェア・通信回線・サービスが組み合わされている
    \item[\kw{POSシステム（Point Of Sales system）}]
      商品を販売した時点でそのデータを収集・管理するシステム。販売管理・在庫管理・生産管理・配送管理などが連携している
    \item[\kw{GPS（Global Positioning System）}]
      複数の衛星の位置情報から現在位置を算出し，バスなどの運行状況通知などに利用できるシステム
    \item[\kw{トレーサビリティシステム}]
      商品の製造番号や電子タグを取り付けることで，生産者から消費者までの商品を追跡可能にする仕組み
  \end{description}
\end{frame}

% --- Slide 3: table ---

\begin{frame}{情報システムの種類}
  \centering
  \begin{tabular}{lll}
    \toprule
    種類 & 説明 & 例 \\
    \midrule
    計算システム & 金銭の出納の集計や気象データの解析など，計算処理を中心とした業務に使われる & 気象予測・統計処理 \\
    通信システム & 気象データの収集や遠隔地とのテレビ会議など，高速で安定したデータ通信を行う業務に使われる & 気象衛星・遠隔医療 \\
    制御システム & 工場の生産設備を自動運転する業務などに使われる & FA（ファクトリーオートメーション） \\
    \bottomrule
  \end{tabular}
\end{frame}

% --- Slide 4: twocol ---

\begin{frame}{情報システムと通貨}
  \begin{columns}[T]
    \begin{column}{0.48\textwidth}
      \begin{block}{電子マネー}
        \small
        \begin{itemize}
          \item 貨幣価値をデジタルデータで表現したもの
          \item IC型：カードやスマートフォンのICチップに金額データを記録（交通系・流通系）
          \item ネットワーク型：インターネット上で金銭情報を管理するものや，プリペイドカードなど
          \item IC型の電子マネーにはRFID（非接触型）の技術が使われている
        \end{itemize}
      \end{block}
    \end{column}
    \begin{column}{0.48\textwidth}
      \begin{exampleblock}{仮想通貨（暗号資産）}
        \small
        \begin{itemize}
          \item インターネット上で新たに定義されたデジタル化された通貨
          \item 既存の各国の貨幣価値と交換でき，物品・サービスの対価として使用できる
          \item ブロックチェーンという技術で，銀行のような機関を経由せずに貨幣価値を移せる
          \item 公的な発行主体や管理者の裏付けなしに流通しているものがあり，価値が暴落・消失することがある
        \end{itemize}
      \end{exampleblock}
    \end{column}
  \end{columns}
\end{frame}

% --- Slide 5: table ---

\begin{frame}{情報システムの処理形態（3種類）}
  \centering
  \begin{tabular}{lll}
    \toprule
    種別 & 説明 & 特徴 \\
    \midrule
    集中情報システム & 中央に大型・中型の汎用コンピュータを設置し，全てのデータを集中管理する & 機密性・完全性が保たれやすい。データ量増加による負荷集中・故障時のリスクがある \\
    分散情報システム & 複数の小型コンピュータを相互に接続し，データを共有しながら業務を分担して処理する。クライアントサーバシステムが代表例 & 負荷が分散されて高速化が望める。セキュリティが低下したりデータの一貫性を保つのが難しかったりする \\
    移動体情報システム & 携帯情報端末を用い，無線技術によって移動中も相互利用を可能にする & 屋内と屋外で別々の回線を切り替えて運用できる柔軟性がある \\
    \bottomrule
  \end{tabular}
\end{frame}

% --- Slide 6: points ---

\begin{frame}{オープンデータ}
  \begin{block}{ポイント}
    \begin{itemize}
      \item 国や地方公共団体，民間企業が所有しているデータのうち，次の①〜③に当てはまるように公開されたものをオープンデータという
      \item ① 営利目的・非営利目的を問わず二次利用可能なルールが適用されたもの
      \item ② 機械判読に適したもの
      \item ③ 無償で利用できるもの
      \item AEDや消火器の設置場所，地域バスの運行ルートなどを公開することで，それらを活用したサービスやアプリが多数登場している
    \end{itemize}
  \end{block}

  \vfill

  \begin{exampleblock}{補足}
    \begin{itemize}
      \item クリエイティブ・コモンズ（CC）：著作者が使用許諾の条件を積極的に開示することで，著作物を再利用されやすくする仕組み
    \end{itemize}
  \end{exampleblock}
\end{frame}

% --- Slide 7: points ---

\begin{frame}{情報システムの信頼性（1）：バックアップとRAID}
  \begin{block}{ポイント}
    \begin{itemize}
      \item 情報システムは稼働を始めると容易に停止することができない。そのため冗長化（二重化・多重化）といった方策が取られる
      \item バックアップ：データを記録した媒体が破損して失われないよう，一定時間で複製を行う方法
      \item フルバックアップ：全てのデータを複製する
      \item 差分バックアップ：前回のフルバックアップからの変更分を複製する
      \item 増分バックアップ：前回のバックアップからの変更分のみを複製する（容量が少なく済む）
      \item RAID：複数のドライブを組み合わせて仮想的に1台の装置に見立てて運用する技術。耐故障性を高められる
    \end{itemize}
  \end{block}
\end{frame}

% --- Slide 8: points ---

\begin{frame}{情報システムの信頼性（2）：稼働率とフェールセーフ}
  \begin{block}{ポイント}
    \begin{itemize}
      \item 稼働率：情報システムが正常に動作している時間の割合を表す平均的な指標
      \item 平均故障間隔（MTBF）：システムが稼働を開始してから故障するまでの平均時間
      \item 平均修理時間（MTTR）：故障が発生してから修理が完了するまでの平均時間
      \item 稼働率 = MTBF ÷ （MTBF + MTTR）
      \item フェールセーフ：機器が故障したとき，安全側に動作するように設計すること（例：信号が消えたら赤になる）
      \item フールプルーフ：誤操作しても危険が起きないようにする設計（例：電池を逆に入れられない形状）
    \end{itemize}
  \end{block}
\end{frame}

% --- Slide 9: exercise ---

\begin{frame}{演習}
  {\small 各自で取り組んでください。}

  \vfill

  \begin{block}{基本問題}
    \begin{itemize}
      \item 集中情報システムと分散情報システムの違いを，機密性・完全性・可用性の観点から説明せよ
      \item オープンデータの3つの条件を説明し，身近な活用例を1つ挙げよ
      \item MTBFが900時間，MTTRが100時間のシステムの稼働率を求めよ
    \end{itemize}
  \end{block}

  \vfill

  \begin{exampleblock}{応用問題（余裕がある人）}
    \begin{itemize}
      \item フェールセーフとフールプルーフの違いを説明し，機械工学の分野での応用例をそれぞれ1つ考えよ
    \end{itemize}
  \end{exampleblock}
\end{frame}

% --- Slide 10: assignment ---

\begin{frame}{課題・次回予告}
  \begin{importantbox}[課題（次回までに提出）]
    教科書 pp.112-117 を復習すること。次週はデータベースの基礎を学ぶ（pp.118-121）
  \end{importantbox}

  \vfill

  \begin{block}{次回}
    「\NextTopic 」（教科書 \NextPages ）
  \end{block}

  \vfill

  {\small
  質問があれば授業後またはTeamsで受け付けます。}
\end{frame}

% --- チェックリスト ---

\begin{frame}{まとめ・チェックリスト}
  \begin{block}{自己チェック（できたら $\checkmark$ をつけよう）}
    \begin{itemize}
      \item[$\square$] 情報システムの種類（集中・分散・移動体）と役割を理解した
      \item[$\square$] 電子マネー・オープンデータなど身近な情報システムを学んだ
      \item[$\square$] 稼働率・RAID・バックアップなど信頼性確保の手法を把握した
    \end{itemize}
  \end{block}

  \vfill

  {\small 質問があれば授業後またはTeamsで受け付けます。}
\end{frame}

\end{document}
